\chapter{Conclusioni}

Questa tesi ha affrontato il problema dell'utilizzo di bpftrace per il tracing di applicazioni Android, partendo dalla constatazione che gli strumenti di osservabilità tipici dei sistemi Linux non sono direttamente disponibili su Android senza adattamenti significativi. La soluzione proposta, basata sull'esecuzione di bpftrace all'interno di una distribuzione Debian installata nell'emulatore Cuttlefish, ha dimostrato che è possibile raccogliere dati di sistema a livello kernel su Android senza modificare il sistema operativo e senza richiedere accesso hardware fisico.

Il sistema sviluppato comprende sette probe bpftrace parallele, un modulo di orchestrazione in Python che coordina la raccolta degli eventi, un modulo di analisi post-sessione e un'interfaccia utente bash. Le probe coprono il ciclo di vita dei processi, le syscall di accesso a file e rete, le transazioni IPC Binder e il volume di dati scambiati via socket. L'architettura a pipeline consente di attivare solo il sottoinsieme di probe necessario per lo scenario di interesse, limitando il volume di dati prodotti e l'overhead sul sistema monitorato.

Le sperimentazioni condotte su tre scenari distinti hanno mostrato che i dati raccolti dal sistema sono sufficientemente precisi da permettere ricostruzioni dettagliate del comportamento delle applicazioni a livello di sistema operativo. Nella sessione di navigazione web con Chrome è stato possibile osservare la catena completa che collega l'interazione dell'utente al traffico di rete: la risoluzione DNS avviata dall'applicazione di ricerca Google, il passaggio dei dati HTTP dal \textit{NetworkService} al renderer Chromium via Binder, e i picchi di attività correlati a ciascuna fase della navigazione. Il volume di traffico IPC interno, circa 17~MB, ha superato di cinque volte il traffico di rete esterno, un dato che da solo non sarebbe emerso dall'analisi di una singola probe.

Nella sessione galleria e fotocamera, il sistema ha permesso di identificare la firma Binder specifica dell'accesso all'hardware fotografico: l'attivazione del thread \texttt{Cam0\_ResultDisp} e del canale verso il Camera HAL è risultata assente durante la sola navigazione della galleria e presente non appena la fotocamera è stata avviata. Questa firma è stabile e riproducibile, e non dipende da indicatori a livello applicativo.

Nella sessione della modalità aereo, la finestra di silenzio nelle probe di associazione socket e accettazione connessioni ha documentato con precisione la disattivazione dello stack WiFi, e la sequenza di riconnessione successiva è risultata ricostruibile passo per passo attraverso i bind di \texttt{wpa\_supplicant}, la negoziazione DHCP di \texttt{IpClient.wlan0} e la creazione dei thread DNS di \texttt{netd}. Il picco Binder che precede di due secondi la ricomparsa dei primi bind di rete ha evidenziato come il \textit{system server} propaghi i cambiamenti di stato della rete a tutti i processi del sistema prima che il collegamento fisico sia effettivamente ristabilito.

In tutti e tre gli scenari, le correlazioni più significative sono emerse dal confronto tra probe diverse, non dall'analisi di una singola fonte di dati. Questo risultato conferma la premessa del lavoro: un sistema di monitoraggio parallelo e multidimensionale permette di osservare il comportamento di Android in modo qualitatively diverso rispetto a strumenti che operano su una sola categoria di eventi.

I limiti del sistema, discussi nel capitolo precedente, riguardano principalmente la dipendenza dall'ambiente emulato e l'assenza di analisi in tempo reale. Entrambi sono limiti architetturali che non invalidano il contributo della tesi, ma definiscono il confine dell'applicabilità attuale: il sistema è uno strumento di analisi e ricerca, non un agente di monitoraggio distribuibile su dispositivi di produzione nella forma attuale.

Il contributo principale di questo lavoro è dimostrare che bpftrace, con l'infrastruttura appropriata, può essere utilizzato per ottenere una visibilità profonda sul comportamento di Android a livello kernel. I dati raccolti hanno permesso di osservare meccanismi interni del sistema operativo, come l'architettura di isolamento del renderer Chromium, la pipeline di elaborazione dei frame della fotocamera e la sequenza di riconnessione WiFi, che non sarebbero accessibili attraverso le API applicative o gli strumenti di profilazione standard. Questo apre prospettive concrete per l'utilizzo di tecniche eBPF nell'analisi della sicurezza delle applicazioni Android, nella verifica del comportamento di componenti di sistema e nel rilevamento di anomalie basato su firme comportamentali a livello kernel.
